\documentclass[12pt,a4paper]{article}
\usepackage[utf8]{inputenc}
\usepackage[T1]{fontenc}
\usepackage[polish]{babel}
\usepackage{geometry}
\usepackage{graphicx}
\usepackage{booktabs}
\usepackage{longtable}
\usepackage{array}
\usepackage{xcolor}
\usepackage{listings}
\usepackage{hyperref}
\usepackage{fancyhdr}
\usepackage{titlesec}
\usepackage{tocloft}
\usepackage{float}

% Ustawienia strony
\geometry{margin=2.5cm}

% Kolory
\definecolor{pgblue}{RGB}{0,83,159}
\definecolor{codegray}{RGB}{245,245,245}
\definecolor{codegreen}{RGB}{0,128,0}

% Styl kodu
\lstset{
    backgroundcolor=\color{codegray},
    basicstyle=\ttfamily\small,
    breaklines=true,
    frame=single,
    numbers=left,
    numberstyle=\tiny\color{gray},
    keywordstyle=\color{pgblue},
    commentstyle=\color{codegreen},
    stringstyle=\color{red},
}

% Nagłówek i stopka
\pagestyle{fancy}
\fancyhf{}
\fancyhead[L]{\textcolor{pgblue}{Dog FACS Dataset - Dokumentacja Techniczna}}
\fancyhead[R]{\textcolor{pgblue}{v2.00}}
\fancyfoot[C]{\thepage}
\renewcommand{\headrulewidth}{0.4pt}

% Formatowanie tytułów
\titleformat{\section}{\Large\bfseries\color{pgblue}}{\thesection}{1em}{}
\titleformat{\subsection}{\large\bfseries\color{pgblue}}{\thesubsection}{1em}{}
\titleformat{\subsubsection}{\normalsize\bfseries}{\thesubsubsection}{1em}{}

\begin{document}

% Strona tytułowa
\begin{titlepage}
    \centering
    \vspace*{2cm}

    {\Large\textbf{Politechnika Gdańska}}\\[0.3cm]
    {\large Wydział Elektroniki, Telekomunikacji i Informatyki}\\[0.3cm]
    {\large Katedra Systemów Decyzyjnych i Robotyki}\\[2cm]

    \rule{\textwidth}{1pt}\\[0.5cm]
    {\Huge\textbf{\textcolor{pgblue}{Dog FACS Dataset}}}\\[0.3cm]
    {\LARGE Dokumentacja Techniczna Projektu}\\[0.5cm]
    \rule{\textwidth}{1pt}\\[2cm]

    {\large\textbf{Pipeline do automatycznej anotacji emocji psów}}\\[0.5cm]
    {\large z wykorzystaniem sztucznej inteligencji}\\[3cm]

    \begin{tabular}{ll}
        \textbf{Wersja:} & 2.00 \\
        \textbf{Data:} & 30 stycznia 2026 \\
        \textbf{Opiekun:} & dr hab. inż. Michał Czubenko \\
    \end{tabular}

    \vfill

    \textbf{Zespół projektowy:}\\[0.3cm]
    \begin{tabular}{lll}
        Danylo Lohachov & 196610 & Kierownik projektu \\
        Anton Shkrebela & 196637 & AI/ML Specialist \\
        Danylo Zherzdiev & 196765 & Backend Developer \\
        Mariia Volkova & 196660 & Data Engineer \\
    \end{tabular}

    \vspace{1cm}
    {\small Rok akademicki 2025/2026}
\end{titlepage}

% Spis treści
\tableofcontents
\newpage

% ============================================================================
\section{Wprowadzenie}
% ============================================================================

\subsection{Cel systemu}

System \textbf{Dog FACS Dataset} to pipeline do automatycznej anotacji emocji psów z wykorzystaniem sztucznej inteligencji. System przetwarza materiały wideo i generuje dataset w formacie COCO zawierający:

\begin{itemize}
    \item \textbf{Bounding boxes} - prostokąty obejmujące wykryte psy
    \item \textbf{Klasyfikację ras} - identyfikacja rasy psa (120+ ras)
    \item \textbf{Punkty kluczowe twarzy} - 20 punktów anatomicznych
    \item \textbf{Action Units} - 11 jednostek akcji mimicznej DogFACS
    \item \textbf{Etykiety emocji} - 6 kategorii (happy, sad, angry, fearful, relaxed, neutral)
\end{itemize}

\subsection{Zakres dokumentu}

Niniejsza dokumentacja techniczna obejmuje:
\begin{enumerate}
    \item Architekturę systemu i pipeline przetwarzania
    \item Specyfikację modeli AI (YOLOv8m, EfficientNet-B4, SimpleBaseline)
    \item System klasyfikacji emocji DogFACS (rule-based)
    \item Format wyjściowy COCO z rozszerzeniami
    \item Wymagania systemowe i stack technologiczny
    \item Interfejs programistyczny (API)
\end{enumerate}

% ============================================================================
\section{Architektura systemu}
% ============================================================================

\subsection{Schemat blokowy}

System składa się z czterech głównych komponentów przetwarzających dane sekwencyjnie:

\begin{verbatim}
    ┌─────────────────────────────────────────────────────────────┐
    │                    Dog FACS Dataset Generator                │
    ├─────────────────────────────────────────────────────────────┤
    │                                                              │
    │  Input (Video/Image)                                         │
    │         │                                                    │
    │         ▼                                                    │
    │  ┌──────────────┐                                           │
    │  │  YOLOv8m     │ ──► Detekcja psów (bounding boxes)        │
    │  │  BBox        │                                           │
    │  └──────┬───────┘                                           │
    │         │                                                    │
    │         ▼ (per detected dog)                                │
    │  ┌──────────────┐    ┌───────────────────┐                  │
    │  │ EfficientNet │    │  SimpleBaseline   │                  │
    │  │ Breed Class. │    │  Keypoints (20)   │                  │
    │  └──────────────┘    └─────────┬─────────┘                  │
    │                                │                             │
    │                                ▼                             │
    │                    ┌───────────────────────┐                │
    │                    │  DogFACS Rule Engine  │                │
    │                    │  (Action Units → Emotion)              │
    │                    └───────────┬───────────┘                │
    │                                │                             │
    │                                ▼                             │
    │                    ┌───────────────────────┐                │
    │                    │   COCO JSON Export    │                │
    │                    └───────────────────────┘                │
    └─────────────────────────────────────────────────────────────┘
\end{verbatim}

\subsection{Architektura warstw}

\begin{table}[H]
\centering
\caption{Warstwy architektury systemu}
\begin{tabular}{|l|l|l|}
\hline
\textbf{Warstwa} & \textbf{Komponenty} & \textbf{Ścieżka} \\
\hline
Application & Streamlit Demo, React Webapp & \texttt{apps/} \\
\hline
Pipeline & InferencePipeline, VideoProcessor & \texttt{packages/pipeline/} \\
\hline
Models & BBox, Breed, Keypoints, Emotion & \texttt{packages/models/} \\
\hline
Data & COCO Handler, Schemas & \texttt{packages/data/} \\
\hline
\end{tabular}
\end{table}

% ============================================================================
\section{Specyfikacja modeli AI}
% ============================================================================

\subsection{Model detekcji psów (BBox)}

\begin{table}[H]
\centering
\caption{Specyfikacja modelu YOLOv8m}
\begin{tabular}{|l|l|}
\hline
\textbf{Parametr} & \textbf{Wartość} \\
\hline
Architektura & YOLOv8m (medium) \\
Plik wag & \texttt{models/yolov8m.pt} \\
Rozmiar pliku & 52.1 MB \\
Rozmiar wejściowy & 640×640 px \\
Wyjście & Lista bounding boxes + confidence \\
Klasy & Filtrowanie tylko klasy "dog" \\
Confidence threshold & 0.3 (domyślnie) \\
\hline
\end{tabular}
\end{table}

\subsection{Model klasyfikacji ras (Breed)}

\begin{table}[H]
\centering
\caption{Specyfikacja modelu EfficientNet-B4}
\begin{tabular}{|l|l|}
\hline
\textbf{Parametr} & \textbf{Wartość} \\
\hline
Architektura & EfficientNet-B4 \\
Plik wag & \texttt{models/breed.pt} \\
Rozmiar pliku & 71.8 MB \\
Rozmiar wejściowy & 224×224 px \\
Wyjście & Top-5 ras z prawdopodobieństwami \\
Liczba klas & 120+ ras \\
Baza wag & ImageNet pretrained \\
\hline
\end{tabular}
\end{table}

\subsection{Model detekcji keypoints}

\begin{table}[H]
\centering
\caption{Specyfikacja modelu SimpleBaseline}
\begin{tabular}{|l|l|}
\hline
\textbf{Parametr} & \textbf{Wartość} \\
\hline
Architektura & SimpleBaseline (ResNet34 + Deconv Head) \\
Plik wag & \texttt{models/keypoints\_dogflw.pt} \\
Rozmiar pliku & 102.1 MB \\
Rozmiar wejściowy & 256×256 px \\
Keypoints natywnie & 46 punktów (DogFLW) \\
Keypoints projektu & 20 punktów (po mapowaniu) \\
\hline
\end{tabular}
\end{table}

\subsubsection{Lista 20 keypoints projektu}

\begin{longtable}{|c|l|l|}
\hline
\textbf{ID} & \textbf{Nazwa} & \textbf{Opis} \\
\hline
\endfirsthead
\hline
\textbf{ID} & \textbf{Nazwa} & \textbf{Opis} \\
\hline
\endhead
0 & left\_eye & Środek lewego oka \\
1 & right\_eye & Środek prawego oka \\
2 & nose & Czubek nosa \\
3 & left\_ear\_base & Podstawa lewego ucha \\
4 & right\_ear\_base & Podstawa prawego ucha \\
5 & left\_ear\_tip & Czubek lewego ucha \\
6 & right\_ear\_tip & Czubek prawego ucha \\
7 & left\_mouth\_corner & Lewy kącik ust \\
8 & right\_mouth\_corner & Prawy kącik ust \\
9 & upper\_lip & Środek górnej wargi \\
10 & lower\_lip & Środek dolnej wargi \\
11 & chin & Podbródek \\
12 & left\_cheek & Lewy policzek \\
13 & right\_cheek & Prawy policzek \\
14 & forehead & Środek czoła \\
15 & left\_eyebrow & Lewa brew \\
16 & right\_eyebrow & Prawa brew \\
17 & muzzle\_top & Góra pyska \\
18 & muzzle\_left & Lewa strona pyska \\
19 & muzzle\_right & Prawa strona pyska \\
\hline
\end{longtable}

% ============================================================================
\section{System DogFACS}
% ============================================================================

\subsection{Action Units (AU)}

System implementuje 11 oficjalnych kodów DogFACS:

\begin{table}[H]
\centering
\caption{Lista Action Units}
\begin{tabular}{|l|l|l|}
\hline
\textbf{Kod AU} & \textbf{Nazwa} & \textbf{Opis} \\
\hline
AU101 & Inner Brow Raiser & Podniesienie wewnętrznej brwi \\
AU102 & Outer Brow Raiser & Podniesienie zewnętrznej brwi \\
AU12 & Lip Corner Puller & Pociągnięcie kącików ust \\
AU25 & Lips Part & Rozchylenie warg \\
AU26 & Jaw Drop & Opadnięcie szczęki \\
AU145 & Blink & Mruganie \\
EAD102 & Ears Forward & Uszy do przodu \\
EAD103 & Ears Flattener & Uszy spłaszczone \\
EAD104 & Ears Rotator & Rotacja uszu \\
AD19 & Tongue Show & Pokazanie języka \\
AD137 & Nose Lick & Lizanie nosa \\
\hline
\end{tabular}
\end{table}

\subsection{Klasyfikacja emocji (Rule-based)}

System klasyfikuje emocje na podstawie kombinacji Action Units:

\begin{table}[H]
\centering
\caption{Kategorie emocji i ich charakterystyki}
\begin{tabular}{|c|l|l|}
\hline
\textbf{ID} & \textbf{Emocja} & \textbf{Charakterystyczne AU} \\
\hline
0 & Happy (szczęśliwy) & AU12, AU25, EAD102 \\
1 & Sad (smutny) & EAD103, AU101 \\
2 & Angry (zły) & AU25, AU26, zmarszczone czoło \\
3 & Fearful (przestraszony) & EAD103, AU145 \\
4 & Relaxed (zrelaksowany) & Miękkie oczy, neutralne uszy \\
5 & Neutral (neutralny) & Brak wyraźnych sygnałów \\
\hline
\end{tabular}
\end{table}

\textbf{WAŻNE:} Klasyfikator emocji jest w pełni oparty na regułach (rule-based) i nie wykorzystuje uczenia maszynowego. Zapewnia to 100\% powtarzalność i interpretowalność wyników.

% ============================================================================
\section{Format wyjściowy COCO}
% ============================================================================

\subsection{Struktura JSON}

\begin{lstlisting}[language=Python, caption=Przykład formatu COCO]
{
  "info": {
    "description": "Dog FACS Dataset",
    "version": "1.0",
    "year": 2025,
    "contributor": "Politechnika Gdanska"
  },
  "images": [
    {
      "id": 1,
      "file_name": "video001_frame_00150.jpg",
      "width": 1920,
      "height": 1080,
      "source_video": "video001.mp4",
      "frame_number": 150
    }
  ],
  "annotations": [
    {
      "id": 1,
      "image_id": 1,
      "category_id": 1,
      "bbox": [100, 150, 300, 400],
      "keypoints": [120.5, 180.3, 2, ...],
      "num_keypoints": 20,
      "breed": "golden_retriever",
      "breed_confidence": 0.95,
      "emotion": "happy",
      "emotion_confidence": 0.87,
      "action_units": {
        "AU12": 1.32,
        "EAD102": 1.18
      }
    }
  ],
  "categories": [
    {
      "id": 1,
      "name": "dog",
      "keypoints": ["left_eye", "right_eye", ...],
      "skeleton": [[0, 1], [0, 2], ...]
    }
  ]
}
\end{lstlisting}

\subsection{Rozszerzenia DogFACS}

Format COCO został rozszerzony o następujące pola specyficzne dla projektu:

\begin{itemize}
    \item \texttt{breed} - nazwa rasy psa
    \item \texttt{breed\_confidence} - pewność klasyfikacji rasy
    \item \texttt{emotion} - etykieta emocji (6 kategorii)
    \item \texttt{emotion\_confidence} - pewność klasyfikacji emocji
    \item \texttt{action\_units} - słownik aktywnych AU z wartościami
    \item \texttt{source\_video} - źródłowe wideo (w images)
    \item \texttt{frame\_number} - numer klatki (w images)
\end{itemize}

% ============================================================================
\section{Wymagania systemowe}
% ============================================================================

\subsection{Minimalne wymagania}

\begin{table}[H]
\centering
\caption{Wymagania systemowe}
\begin{tabular}{|l|l|l|}
\hline
\textbf{Komponent} & \textbf{Minimalne} & \textbf{Zalecane} \\
\hline
System operacyjny & Windows 10, Linux, macOS & Windows 11, Ubuntu 22.04 \\
Python & 3.10+ & 3.11+ \\
RAM & 8 GB & 16 GB \\
Dysk & 15 GB & 50 GB \\
GPU & - & NVIDIA CUDA 11.8+ \\
VRAM & - & 8 GB+ \\
FFmpeg & Wymagany & - \\
\hline
\end{tabular}
\end{table}

\subsection{Stack technologiczny}

\begin{table}[H]
\centering
\caption{Technologie i wersje}
\begin{tabular}{|l|l|l|}
\hline
\textbf{Kategoria} & \textbf{Technologia} & \textbf{Wersja} \\
\hline
Runtime & Python & 3.10+ \\
ML Framework & PyTorch & 2.0+ \\
Detection & Ultralytics (YOLOv8) & 8.0+ \\
Classification & timm (EfficientNet) & 0.9+ \\
Computer Vision & OpenCV & 4.8+ \\
Frontend Demo & Streamlit & 1.28+ \\
Backend API & FastAPI & 0.100+ \\
Linting & Ruff & 0.1+ \\
\hline
\end{tabular}
\end{table}

% ============================================================================
\section{Interfejs programistyczny (API)}
% ============================================================================

\subsection{Pipeline podstawowy}

\begin{lstlisting}[language=Python, caption=Przykład użycia API]
from packages.pipeline import InferencePipeline, PipelineConfig

# Konfiguracja
config = PipelineConfig(
    bbox_weights="models/yolov8m.pt",
    breed_weights="models/breed.pt",
    keypoints_weights="models/keypoints_dogflw.pt",
    device="cuda",
    confidence_threshold=0.3
)

# Inicjalizacja
pipeline = InferencePipeline(config)

# Przetwarzanie obrazu
result = pipeline.process_image(image)

for annotation in result:
    print(f"Breed: {annotation.breed}")
    print(f"Emotion: {annotation.emotion}")
    print(f"Keypoints: {len(annotation.keypoints)}")
\end{lstlisting}

\subsection{COCO Dataset API}

\begin{lstlisting}[language=Python, caption=Tworzenie datasetu COCO]
from packages.data.coco import COCODataset

# Inicjalizacja
dataset = COCODataset()

# Dodanie obrazu
image_id = dataset.add_image(
    file_name="frame.jpg",
    width=1920,
    height=1080,
    source_video="video.mp4",
    frame_number=150
)

# Dodanie anotacji
dataset.add_annotation(annotation, image_id)

# Zapis
dataset.save("output.json")
\end{lstlisting}

% ============================================================================
\section{Repozytorium}
% ============================================================================

\begin{table}[H]
\centering
\caption{Informacje o repozytorium}
\begin{tabular}{|l|l|}
\hline
\textbf{Pole} & \textbf{Wartość} \\
\hline
URL & \url{https://github.com/eternaki/group-project} \\
Gałąź główna & main \\
Gałąź feature & feature/dogfacs-rule-based \\
Struktura gałęzi & main → develop → sprint-X \\
Model storage & Git LFS \\
Licencja & MIT \\
\hline
\end{tabular}
\end{table}

% ============================================================================
\section*{Historia dokumentu}
% ============================================================================

\begin{table}[H]
\centering
\begin{tabular}{|l|l|l|p{6cm}|}
\hline
\textbf{Wersja} & \textbf{Data} & \textbf{Autor} & \textbf{Opis zmian} \\
\hline
0.1 & 03.01.2025 & D. Zherzdiev & Wstępna architektura systemu, definicja pipeline \\
\hline
0.2 & 10.01.2025 & A. Shkrebela & Specyfikacja modeli AI (YOLOv8, EfficientNet) \\
\hline
0.3 & 16.01.2025 & A. Shkrebela & Dodanie systemu DogFACS, keypoints mapping \\
\hline
0.4 & 20.01.2025 & D. Zherzdiev & Specyfikacja formatu COCO z rozszerzeniami \\
\hline
0.5 & 24.01.2025 & M. Volkova & Wymagania systemowe, stack technologiczny \\
\hline
1.0 & 28.01.2025 & D. Lohachov & Konsolidacja dokumentacji, wersja do recenzji \\
\hline
1.1 & 29.01.2025 & A. Shkrebela & Aktualizacja rule-based emotion classifier \\
\hline
2.0 & 30.01.2025 & D. Lohachov & Migracja do formatu LaTeX, finalna wersja \\
\hline
\end{tabular}
\end{table}

\subsection*{Autorzy sekcji}
\begin{itemize}
    \item \textbf{Architektura systemu:} Danylo Zherzdiev
    \item \textbf{Modele AI (BBox, Breed, Keypoints):} Anton Shkrebela
    \item \textbf{System DogFACS:} Anton Shkrebela
    \item \textbf{Format COCO:} Danylo Zherzdiev
    \item \textbf{Wymagania systemowe:} Mariia Volkova
    \item \textbf{API i przykłady:} Danylo Zherzdiev
    \item \textbf{Redakcja i koordynacja:} Danylo Lohachov
\end{itemize}

\vfill
\begin{center}
\textit{Dokumentacja Techniczna Projektu - Dog FACS Dataset}\\
\textit{Politechnika Gdańska, WETI}\\
\textit{Katedra Systemów Decyzyjnych i Robotyki}\\
\textit{Rok akademicki 2025/2026}
\end{center}

\end{document}
